
\chapter{Background and Related Work} \label{CH2_LiteratureReview}

\normalsize\doublespacing

This chapter surveys the landscape of energy-performance co-optimization in high-performance
computing (HPC) at the node level, covering the physical limits that motivate this research,
the theoretical scheduling foundations that frame the optimization problem, the hardware
and software mechanisms available to actuate power, the control and learning strategies
proposed to exploit them, the performance metrics used to close the feedback loop, and
the edge computing platforms that connect scientific instruments to HPC systems.

%% =========================================================
\section{The Exascale Power Wall and Heterogeneous Architectures}
%% =========================================================

The exponential growth in transistor density predicted by Moore's Law delivered roughly
a doubling of single-thread performance every two years for three decades.
Dennard scaling~\cite{Dennard1974}, the companion observation that power density remains
constant as transistors shrink, was the physical foundation of that trend.
As feature sizes fell below~\SI{65}{\nano\meter}, however, leakage current grew super-linearly
and Dennard scaling broke down: frequency stopped climbing, leaving power as the primary
constraint on compute density.

The result is the \emph{power wall}.  Modern leadership-class systems approach or exceed
100\,MW of facility power, and the cost of that electricity over a machine's lifetime can
rival its capital cost.  Procurement agencies at national laboratories and supercomputing
centers now impose hard power-budget constraints at installation time, shifting the design
objective from peak floating-point throughput to performance-per-watt.
Rountree et al.~\cite{rountree2012beyond} showed that even hardware-enforced power bounds
are more nuanced than simple frequency throttling: the relationship between a power cap, the
resulting frequency, and application throughput is non-linear and workload-dependent.

The transition to heterogeneous nodes - pairing multi-core CPUs with one or more
many-core accelerators (GPUs, FPGAs, or co-processors) - compounds the challenge.
Heterogeneous architectures expose a wider dynamic range of power consumption per
device and create multiple distinct power domains that must be controlled jointly.
Static assignment of power to each device ignores the fact that workload phases
shift the bottleneck between the CPU and the accelerator over time.
The analysis by Wang et al.~\cite{10.1145/3716872} demonstrates this point: on
CPU-GPU platforms, decentralized per-device RL agents with centralized coordination
improve performance by 8.5\,\% over state-of-the-art heuristics precisely because
they can track the shifting power demand dynamically.
Maximizing efficiency on such systems therefore requires fine-grained, adaptive
control strategies rather than coarse-grained static allocations.

%% =========================================================
\section{Energy-Aware Scheduling Theory}
%% =========================================================

A theoretical foundation for energy-efficient computing comes from classical
scheduling theory, which provides provably optimal or competitive algorithms for
assigning execution speeds to jobs.
The canonical formulation places a set of jobs with release times, deadlines, and
known workloads on a variable-speed processor and seeks to minimize total energy
consumption while satisfying performance constraints.

The Yao-Demers-Shenker (YDS) algorithm~\cite{YDS} solves this problem optimally
in the offline setting, where all job parameters are known in advance.
YDS identifies intervals of maximum workload density and assigns constant execution
speeds within those intervals, achieving the minimum possible energy while meeting
all deadlines.
Earliest Deadline First (EDF) scheduling is commonly paired with YDS-derived speed
functions as an online policy~\cite{YDS}.
For scenarios without explicit deadlines, prior work considers minimizing a
combination of energy and flow time, yielding different optimality criteria and
competitive guarantees~\cite{ALBERS20151194}.
Online algorithms are evaluated using competitive analysis: an algorithm is
\emph{c-competitive} if its objective value is within a factor $c$ of the optimal
offline solution for any input~\cite{YDS,ALBERS20151194}.

While these formulations provide strong theoretical insights, they assume precise
knowledge of job characteristics and speed-power relationships.
Modern HPC applications exhibit complex, hardware-dependent behavior that
invalidates closed-form speed-power models, and workloads arrive dynamically
without advance knowledge of future jobs.
This gap between scheduling theory and practice motivates the data-driven,
application-agnostic control approaches presented in Chapters~\ref{CH3_IC2E}
through~\ref{CH5_ICS}.

%% =========================================================
\section{Hardware Power Actuators}
%% =========================================================

Effective software power management depends on hardware mechanisms that expose
the power-performance trade-off as a controllable degree of freedom.
This section reviews the principal actuators available on contemporary HPC nodes.

\subsection{Dynamic Voltage and Frequency Scaling}

Dynamic voltage and frequency scaling (DVFS) has been the canonical hardware
power actuator for three decades.
Because dynamic power scales as $P \propto C V^2 f$, reducing the operating
frequency~$f$ (and the supply voltage~$V$ that must accompany it to maintain
timing margins) yields a cubic reduction in dynamic power at the cost of longer
execution time.
The Linux \texttt{cpufreq} framework~\cite{wysocki2017cpufreq} exposes
P-state selection to the operating system and user-space governors, enabling
coarse-grained software control.
Poellabauer et al.~\cite{DVFS_Control} demonstrated feedback-driven DVFS for
memory-bound real-time applications: by using cache-miss rates as a leading
indicator of future CPU demand, the controller reduces energy while meeting
timing constraints more reliably than utilization-only predictors.
Chen et al.~\cite{DVFS_optimization} extended DVFS to multiprocessor
system-on-chip designs, combining it with dynamic power management (DPM)
sleep states to find energy-optimal configurations for real-time task sets.

A critical limitation of DVFS on modern CPUs is that its effectiveness is
workload-dependent.  Compute-bound applications that fully utilize the
execution pipeline gain little energy benefit from frequency reduction - the
longer execution time offsets the lower instantaneous power.
Cochran et al.~\cite{cochran2011pack} addressed this by coupling DVFS with
thread packing: consolidating threads onto fewer cores allows unused cores to
enter deep sleep, recovering static leakage savings that DVFS alone cannot
achieve.

\subsection{Intel Running Average Power Limit}

Intel's Running Average Power Limit (RAPL)~\cite{david2010rapl} interface
introduced hardware-enforced power capping directly in the processor and memory
controller.  Unlike DVFS, which adjusts voltage and frequency in discrete
P-state steps, RAPL enforces a time-averaged power constraint by throttling
clock delivery across an entire power domain - the package, the DRAM, or
individual cores - on a millisecond timescale.
This makes RAPL a powerful actuator for software controllers because it
abstracts away the detailed mechanics of frequency selection and exposes a
single power-budget knob per domain.

RAPL's package power limit (PKG) controls processor-side consumption, while
the DRAM power limit controls memory-side consumption independently.
For memory-bandwidth-bound applications, capping DRAM power can shift the
bottleneck back to the CPU, enabling higher CPU frequency at the same node
power budget.
Petoumenos et al.~\cite{petoumenos2015power} characterized the conditions under
which RAPL provides energy savings and showed that the relationship between the
RAPL cap value and achieved power is non-trivial: power may exceed the cap
transiently during burst phases and fall below it during idle phases.

\subsection{Uncore Frequency Scaling}

Beyond the compute cores, modern Intel processors expose a separate frequency
domain for the uncore - the last-level cache, the memory controllers, and the
ring or mesh interconnect.  Uncore frequency scaling (UFS) allows the operating
system or a runtime to reduce this domain's clock independently of the core
frequency, yielding savings for memory-bound workloads that would otherwise idle
the uncore at its peak frequency between bursts.
The Intel Speed Select Technology and \texttt{intel\_uncore\_frequency} Linux
driver make UFS accessible to system software.
Studies have shown UFS savings of up to 13\,\% in node-level energy for
memory-bandwidth-bound workloads, making it an important complementary actuator
to RAPL and core DVFS.

\subsection{Processor Idle States}

Modern processors expose a hierarchy of sleep states (C-states) that trade
wake-up latency for progressively deeper power reduction.  C0 is the active
state; C1 halts the core clock; C6 powers down the core voltage rail
entirely, saving the largest fraction of static leakage at the cost of a
microsecond-scale wake-up penalty.  Package C-states extend this hierarchy to
encompass the entire socket: if all cores are in a sufficiently deep C-state,
the package can power-gate shared structures such as the LLC and memory
controllers.  Intelligent use of idle states requires that the OS scheduler,
power governor, and application runtime cooperate to avoid excessive wake-up
overhead while capturing the maximum savings during idle intervals.

%% =========================================================
\section{Software Runtime Strategies for Power Management}
%% =========================================================

Hardware actuators provide the means; software runtimes provide the policy.
The literature spans three broad categories: heuristic and profile-driven
runtimes, control-theoretic feedback controllers, and data-driven machine
learning approaches.

\subsection{Heuristic and Profile-Driven Runtimes}

The Energy Aware Runtime (EAR)~\cite{corbalan2019ear} is a production-grade
runtime deployed on Tier-1 European systems.  EAR detects MPI communication
phases by monitoring MPI call frequency, then reduces the RAPL power cap or the
core frequency during identified communication-idle windows when the application
would otherwise burn power while stalled on collective operations.
The key insight is that MPI barriers and collectives create predictable
performance-insensitive intervals during which power can be reduced without
affecting time-to-solution.

A complementary approach is taken by COUNTDOWN~\cite{countdown2021libri}, which
intercepts MPI calls via a preloaded shared library and reduces core frequency
during the waiting time within collectives.  By operating transparently at the
MPI level, COUNTDOWN requires no application modification and has been validated
across production workloads on European HPC systems.

The Global Extensible Open Power Manager
(GEOPM)~\cite{eastep2017global} takes a hierarchical approach, coordinating
power across nodes via a tree of runtime agents.  A per-node agent observes
application progress and reallocates its fraction of the system-wide power
budget in response.  On CORAL procurement benchmarks on a Xeon~Phi cluster,
GEOPM's rebalancing plugin improved time-to-solution by up to 30\,\% over
a static power allocation.
The Bull Dynamic Power Optimizer~\cite{8514915} similarly uses in-band
performance counters - specifically instructions-per-cycle - to
dynamically modulate core frequencies during execution.

Patki et al.~\cite{patki2015practical} demonstrated that power-constrained
scheduling can improve job throughput by over-provisioning nodes at reduced
frequency, fitting more concurrent jobs within a facility power budget.
The Systemwide Power Management framework built on the Argo NodeOS~\cite{Ellsworth:16:Systemwide}
extended this to the full system stack, linking job-scheduler power allocations
with per-node runtime actuators.

Profile-driven approaches such as the Dynamic Energy-Performance
Optimizer (DEPO)~\cite{KRZYWANIAK2023396} pre-collect per-phase power
profiles and replay them at runtime to match observed phases.
While effective when training and deployment workloads are identical,
such approaches degrade when workload characteristics shift, motivating
adaptive and learning-based alternatives.

A hybrid hardware-software approach is taken by the Performance under Power Limits
(PUPiL) framework~\cite{zhang2016maximizing}, which jointly manipulates DVFS, core
allocation, socket usage, memory usage, and hyperthreading to achieve energy and
power consumption benefits beyond what raw RAPL capping alone can deliver.
Score-P~\cite{czarnul2019energy}, designed for analysis and optimization of HPC
applications, further enables energy-aware analysis and distinguishes between
energy-optimal and time-optimal configurations.
The effects of DVFS on both CPU and GPU components have been characterized
by Ge et al.~\cite{ge2013effects}, who showed that controlling core and uncore
frequencies independently improves efficiency across heterogeneous compute nodes.

\subsection{Power Measurement and Profiling}

Before any runtime can regulate power, it must first quantify what the system is
consuming.
Ryan et al.~\cite{8323587} surveyed methods for analyzing application-level power
profiles at runtime, classifying them as \emph{in-band} or \emph{out-of-band}
based on the measurement approach.
Out-of-band methods use external meters to obtain job-level energy totals during
or after execution; these are accurate but provide no fine-grained visibility into
transient power phases.
In-band methods read device-integrated sensors, such as Intel RAPL counters, at
millisecond timescales, enabling dynamic power management by exposing instantaneous
consumption to the controlling software.
Raghavendra et al.~\cite{raghavendra2008no} demonstrated early use of in-band
sensors to cut power to subsystems by detecting no-load zones during application
execution, achieving fine-grained power regulation in a hardware-centric manner.
Simmendinger et al.~\cite{10740868} later demonstrated that dynamic power capping,
rather than fixed power limits, is essential for handling behavioral changes in
applications and for redistributing available power across nodes or resources
within a datacenter.
Together, these profiling and measurement capabilities form the observability layer
on which all runtime power controllers depend.

\subsection{Control-Theoretic Feedback Controllers}

Control theory provides a principled framework for designing power managers
that track a performance set-point despite disturbances and model uncertainty.
Cerf et al.~\cite{cerf2021sustaining} designed a classical
proportional-integral-derivative (PID) controller that adjusts the RAPL power
cap to maintain a user-specified performance level measured via hardware
performance counters.  The controller was demonstrated to reject step
disturbances in workload intensity while converging to within a few percent of
the target performance.

Hawila et al.~\cite{ccta22} extended this to adaptive control using a
model-reference architecture.  Rather than tuning PID gains manually, the
adaptive controller estimates the plant model (the mapping from RAPL cap to
application progress) online and adjusts gains to track the reference model's
response.  Evaluated on several HPC benchmarks, the approach achieved stable
performance tracking without workload-specific re-tuning.

CoPPer~\cite{CoPPer} combines hardware power capping with a feedback loop
that targets soft real-time performance goals expressed as heartbeats.
The controller modulates the RAPL power limit to track an application-specified
throughput target, automatically trading between energy savings and performance
based on a configurable preference.
Ramesh et al.~\cite{ramesh2019understanding} quantified the impact of RAPL
throttling on application progress, showing that the phase-dependent response
of applications to power caps must be accounted for in controller design.

A shared limitation of control-theoretic approaches is that they require
a reliable performance signal to close the feedback loop.  When the available
metrics fail to capture true application progress - for example, when IPC
conflates useful work with stall cycles - the controller may incorrectly diagnose
performance degradation and respond counterproductively.
This motivates the careful study of performance proxies discussed in
Section~\ref{sec:litrev_perf_proxies}.

\subsection{Machine Learning and Reinforcement Learning Approaches}

Machine learning replaces hand-tuned heuristics and plant models with
data-driven policies that can generalize across workload conditions.
Multi-objective machine learning (MOML) methods~\cite{MOML} construct
performance and energy models from micro-benchmark traces and then identify
Pareto-optimal power-cap configurations, exposing the energy-delay
trade-off as an explicit frontier.

Reinforcement learning (RL) frames power management as a sequential
decision-making problem: at each time step the controller observes the system
state (e.g., current power, performance, and application phase), takes an action
(a power-cap value), and receives a reward that encodes the desired
energy-performance trade-off.  The controller learns a policy that maximizes
cumulative reward without requiring an explicit system model.
Wang et al.~\cite{10.1145/3716872} proposed a multi-agent deep RL (MADRL)
framework for CPU-GPU heterogeneous platforms, assigning an independent RL agent
to each device.  A centralized coordinator maximizes utilization within the total
power budget by dynamically reallocating power between agents, yielding an
average performance improvement of 8.5\,\% over heuristic baselines.

Early learning-based methods applied Q-learning and Double-Q learning
(DDQL)~\cite{DDQL} to DVFS governors, enabling systems to adapt frequency
decisions to observed workload behavior without an explicit system model.
These online approaches improve energy savings over static governors, but their
convergence time and sensitivity to hyperparameters limit their applicability in
production HPC environments where exploration can interfere with running jobs.
More recent efforts employ model-based or multi-objective learning techniques to
predict performance and energy under different power-cap configurations~\cite{MOML},
constructing performance and energy models from micro-benchmark traces and then
identifying Pareto-optimal configurations across diverse
workloads~\cite{HPC_proxy_applications}.

Online RL requires live interaction with the production system during training,
which can interfere with running applications, induce thermal or power
excursions, and incur significant computational overhead.
Offline RL~\cite{levine2020offline} addresses this by learning from a static
dataset of previously collected execution traces, decoupling training from
deployment.  The key algorithmic challenge is distributional shift: the trained
policy may encounter state-action combinations absent from the dataset, leading
to overestimated Q-values and unsafe actions.  Conservative Q-Learning
(CQL)~\cite{levine2020offline} mitigates this by penalizing out-of-distribution
Q-values during training.

Raj et al.~\cite{ic2e23} applied online RL to node-level power capping using
a PPO agent trained on a mathematical model of a STREAM-based compute node,
demonstrating performance-aware power reduction through RAPL actuators.
Building on this, the authors extended the approach to offline RL, learning an
application-agnostic policy from pre-collected execution traces across diverse
benchmarks~\cite{offlineRL} and validating it on out-of-distribution workloads.

Multi-objective RL explicitly models the energy-performance Pareto front.
The Preference-Driven MORL (PDMORL) framework~\cite{basaklar2023pdmorlpreferencedrivenmultiobjectivereinforcement}
trains a single policy conditioned on a user-provided preference vector that
selects a point on the Pareto front at deployment time, eliminating the need
for separate agents for each trade-off configuration.

Recent work has scaled RL-based power management beyond the individual compute
node to the cluster and job-scheduler level.
Amrizal et al.~\cite{amrizal2025spars} introduced SPARS, a simulation environment
for training RL agents that decide when to power down idle cluster nodes, balancing
the energy savings from node shutdown against the wake-up latency penalties that
affect job throughput.
Budiarjo et al.~\cite{budiarjo2025curriculum} further improved the sample efficiency
of deep RL cluster power management by applying curriculum learning, starting the
agent in simplified scenarios and progressively increasing workload complexity, which
accelerates convergence and improves the quality of the learned policy.
These cluster-level approaches are complementary to the node-level RAPL-based
controllers presented in this dissertation: they decide \emph{which} nodes to keep
active, while node-level controllers manage \emph{how much} power an active node
may consume.

%% =========================================================
\section{Performance Proxies for Power Management}
\label{sec:litrev_perf_proxies}
%% =========================================================

Closing a power-management feedback loop requires a runtime performance signal
that is (i) available online without application instrumentation, (ii)
responsive to changes in the power cap on a timescale that matches the
controller's action interval, and (iii) representative of useful application
progress rather than low-level micro-architectural activity.
Selecting an inappropriate proxy leads to control decisions that either waste
energy or degrade performance unnecessarily.

\subsection{Instructions Per Cycle}

Instructions per cycle (IPC) is the most widely used hardware performance
counter for characterizing CPU efficiency, and it is naturally available through
standard interfaces such as PAPI~\cite{mucci1999papi} and
\texttt{perf\_event}.  IPC measures the ratio of retired instructions to elapsed
clock cycles.

However, IPC is an unreliable proxy for application-level progress under power
capping.  Ramesh et al.~\cite{ramesh2019understanding} demonstrated that IPC
can remain stable or even increase when a power cap slows execution, because
reducing the clock frequency shrinks the denominator (elapsed cycles at the new
frequency) proportionally.  More critically, IPC conflates useful compute work
with instruction-level parallelism and out-of-order execution overhead: a core
stalled awaiting a cache miss executes no useful instructions but still
contributes to the instruction count denominator.  For memory-bandwidth-bound
workloads - a dominant class in HPC - IPC therefore fails to distinguish the
performance-limited regime from the power-limited regime.

\subsection{Application Heartbeats}

The heartbeat abstraction~\cite{acar2018heartbeat} decouples the
application-level notion of progress from micro-architectural metrics.
Applications emit a ``heartbeat'' (an API call) at regular semantic
checkpoints - e.g., at the end of each time-step in a simulation, or after each
problem partition is solved.  The heartbeat rate (heartbeats per second) is
directly interpretable as application throughput, independent of architectural
details such as frequency or cache behavior.  A power manager monitoring
heartbeat rate can therefore distinguish a genuine throughput reduction caused by
an overly aggressive power cap from a change in IPC that merely reflects
micro-architectural reorganization.

The limitation of heartbeats is that they require application instrumentation,
which is impractical for third-party or legacy codes in production environments.
This has motivated work on application-agnostic proxies derived purely from
hardware counters.

\subsection{Memory Requests Per Instruction and Related Metrics}

For HPC workloads that are memory-bandwidth-bound, the ratio of last-level cache
misses to retired instructions - sometimes called the Cache Miss Rate per
instruction (CMR) or Memory Requests Per Instruction (MRPI) - captures the
degree to which execution progress is limited by off-chip bandwidth rather than
compute throughput.  When the power cap reduces core frequency, a
memory-bandwidth-bound application's execution time increases in proportion,
while its CMR remains approximately constant: each instruction still generates
the same number of off-chip requests.  This makes CMR/MRPI a more stable proxy
for progress than IPC under DVFS and RAPL throttling.

Several works have leveraged this property.  The Bull Dynamic Power
Optimizer~\cite{8514915} uses IPC as a phase-change indicator in combination
with DVFS, but acknowledges that memory-bound phases require a separate
characterization.  The offline RL approach described in
Chapter~\ref{CH4_HPDC} adopts an application-agnostic progress metric
derived from hardware counters and normalized to the full-power baseline, enabling
comparison across applications with different computational footprints.

\subsection{Summary}

Table~\ref{tab:litrev_proxies} summarizes the performance proxies discussed
above along with their key properties.  The central tension is between
ease of availability (favoring hardware counters) and semantic fidelity
(favoring application-level heartbeats).  The research presented in this
dissertation navigates this tension by adopting a normalized progress metric
derived from hardware performance counters, combining the availability of
IPC-class measurements with improved robustness to frequency changes.

\begin{table}[t]
  \centering
  \caption{Comparison of performance proxies for power management.}
  \label{tab:litrev_proxies}
  \begin{tabular}{lccc}
    \toprule
    \textbf{Proxy} & \textbf{Requires instrumentation} & \textbf{Frequency-robust} & \textbf{Semantic fidelity} \\
    \midrule
    IPC                & No  & No  & Low  \\
    CMR / MRPI         & No  & Yes & Medium \\
    Heartbeats         & Yes & Yes & High \\
    \bottomrule
  \end{tabular}
\end{table}

%% =========================================================
\section{Edge Computing and Edge-to-HPC Workflows}
%% =========================================================

Large-scale scientific experiments increasingly rely on computing systems that span
from field-deployed sensors and instruments to centralized HPC clusters.
In this \emph{edge-to-HPC continuum}, data is generated at remote edge nodes,
processed locally to reduce volume and latency, and forwarded to HPC systems for
full-scale simulation and analysis.
Managing the resources of this continuum requires mechanisms that operate across
heterogeneous hardware and bridge vastly different software stacks.

\subsection{Container-Based Autoscaling}

Modern edge deployments commonly use container orchestration platforms such as
Kubernetes to manage application lifecycle and resource allocation.
Proactive Kubernetes autoscalers~\cite{ju2021proactive} control the number of
application containers and their resource allocations based on predicted demand,
while survey work~\cite{tran2022survey} catalogs the breadth of autoscaling
policies in use across production deployments.
However, these autoscalers operate primarily on application-level signals and do
not consider metrics at the workflow or infrastructure level, limiting their
ability to coordinate resource management across the full edge-to-HPC stack.
Furthermore, these policies are threshold-based and do not accommodate use-case
specific designs, leaving a gap when workload-adaptive control is required.

\subsection{Edge Computing Platforms for Scientific Instruments}

The Waggle platform~\cite{beckman2016waggle} is a field-deployable edge computing
node designed for environmental and scientific data collection and analysis in the
field.
Waggle provides scheduling capabilities that allow multiple users to share node
resources and execute code at the edge~\cite{kim2022goal}.
Early use of these scheduling capabilities revealed the need for finer-grained
resource management to serve multiple concurrent applications with limited
computational and storage budgets.

To address this need, the Argo Node Resource Manager~\cite{argo-ipdps-2019}
enables application-specific resource monitoring at runtime and supports dynamic
resource arbitration: it tracks and reacts to changes in application behavior over
the lifetime of a scientific experiment.
These capabilities make Argo a natural fit for the edge tier of an integrated
scientific computing workflow.
The Systemwide Power Management framework~\cite{Ellsworth:16:Systemwide} similarly
extends resource arbitration to the full system stack, linking job-scheduler power
allocations with per-node runtime actuators.

\subsection{Data Streaming Across the Continuum}

Moving data efficiently between edge nodes and HPC systems is a prerequisite for
any edge-to-HPC workflow.
George et al.~\cite{george2025edge2hpc} characterized three cross-facility data
streaming architectures, namely direct streaming, proxied streaming, and managed
service streaming, analyzing their trade-offs in data flow path, deployment
complexity, and achievable throughput.
Their analysis using the SciStream memory-to-memory streaming framework shows that
architectural choice has a significant impact on end-to-end latency and resource
utilization, and that no single architecture dominates across all use cases.
This motivates the need for adaptive middleware that can select and configure the
appropriate streaming path based on current system conditions and workflow demands.

\subsection{Edge AI Model Optimization Using HPC}

A critical capability in the edge-to-HPC continuum is the ability to train or
fine-tune AI models deployed at the edge using the full computational power of
connected HPC systems.
Aach et al.~\cite{aach2025edgeai} demonstrated a workflow in which edge AI models
for quality control in additive manufacturing are optimized on HPC systems with
the edge device active in the training loop, enabling hardware-aware model
compression that would be impractical to perform on the edge device itself.
The resulting models maintain high accuracy while satisfying the memory and latency
constraints of the target edge hardware.
This edge-in-the-loop training paradigm is representative of the broader class of
workflows that motivate cross-continuum resource management frameworks.

\subsection{Dynamic Resource Arbitration Across the Continuum}

While autoscalers and platform runtimes individually improve resource utilization,
connecting these mechanisms across the continuum remains challenging.
Applications running at the edge need access to HPC systems for model retraining
and large-scale analysis, requiring resource reservation and data movement
coordination that current frameworks do not handle automatically.
AI models deployed at the edge also need to be retrained with new experimental data
using HPC resources, allowing workflows to adapt as scientific needs evolve.
Chapter~\ref{CH6_Charon} presents the Charon framework, which addresses how edge
resources can be arbitrated dynamically under user-provided performance goals, how
system software can support AI-at-the-edge applications that leverage HPC platforms
for model tuning, and how flexible policies can be deployed across diverse hardware
setups and use cases.
